%==============================================================================
% Chapitre 6 - Atmosphère standard
%==============================================================================

\section{Introduction}

Lorsqu'un projectile quitte le canon d'une arme, il ne se déplace pas dans le
vide.

Il évolue au sein de l'atmosphère terrestre, un milieu dont les propriétés
varient continuellement.

La température, la pression et la densité de l'air influencent directement
les performances balistiques.

Afin de disposer d'une référence commune, les ingénieurs et les scientifiques
utilisent un modèle appelé \emph{atmosphère standard}.

Ce modèle permet de comparer des résultats obtenus dans des conditions
différentes et constitue la base de nombreux calculs balistiques.

\begin{aretenir}

L'atmosphère standard est un modèle de référence.

Elle ne décrit pas exactement l'atmosphère réelle à un instant donné, mais
fournit une base commune pour les calculs et les comparaisons.

\end{aretenir}

\section{Pourquoi définir une atmosphère standard ?}

Les conditions météorologiques changent constamment.

Au cours d'une même journée, la température, la pression et l'humidité peuvent
évoluer de manière significative.

Si aucun modèle de référence n'existait, il serait difficile :

\begin{itemize}
\item de comparer des essais ;
\item de construire des tables balistiques ;
\item de développer des modèles numériques ;
\item de comparer les performances de différentes munitions.
\end{itemize}

L'atmosphère standard fournit donc un environnement de référence dans lequel
les calculs peuvent être réalisés.

\section{Atmosphère standard internationale}

Le modèle le plus utilisé est l'\emph{International Standard Atmosphere}
(ISA).

Ce modèle définit notamment :

\begin{itemize}
\item la température ;
\item la pression ;
\item la densité de l'air ;
\item la vitesse du son ;
\end{itemize}

en fonction de l'altitude.

\section{Conditions de référence au niveau de la mer}

Dans le modèle ISA, les conditions de référence au niveau moyen de la mer sont :

\begin{center}
\begin{tabular}{ll}
\toprule
Grandeur & Valeur \\
\midrule
Altitude & 0 m \\
Température & 15°C \\
Pression & 101,325 Pa \\
Densité de l'air & 1,225 kg/m$^3$ \\
Vitesse du son & 340 à 343 m/s \\
\bottomrule
\end{tabular}
\end{center}

Ces valeurs seront fréquemment utilisées dans les chapitres suivants.

\section{Variation avec l'altitude}

L'atmosphère n'est pas homogène.

Lorsque l'altitude augmente :

\begin{itemize}
\item la pression diminue ;
\item la densité diminue ;
\item la température évolue.
\end{itemize}

Pour un projectile, cela signifie que l'air oppose généralement moins de
résistance à haute altitude.

\begin{exemple}

À 2,000 mètres d'altitude, un projectile conserve généralement mieux sa
vitesse qu'au niveau de la mer.

La portée et la trajectoire peuvent alors être sensiblement modifiées.

\end{exemple}

\section{Influence sur la balistique}

Les conditions atmosphériques influencent directement :

\begin{itemize}
\item la traînée ;
\item le temps de vol ;
\item la chute du projectile ;
\item la dérive due au vent ;
\item la vitesse résiduelle ;
\item la portée utile.
\end{itemize}

Deux tirs réalisés avec la même arme et la même munition peuvent donc produire
des résultats différents si les conditions atmosphériques changent.

\section{Altitude-densité}

En aéronautique et en balistique, il est souvent plus utile de considérer
la densité réelle de l'air que l'altitude géométrique.

On parle alors d'\emph{altitude-densité}.

Cette grandeur tient compte simultanément :

\begin{itemize}
\item de l'altitude ;
\item de la température ;
\item de la pression.
\end{itemize}

Deux lieux situés à la même altitude géométrique peuvent présenter des
altitudes-densité très différentes.

\begin{exemple}

Une journée très chaude en montagne peut produire une altitude-densité
supérieure à celle observée lors d'une journée froide à une altitude pourtant
plus élevée.

\end{exemple}

\section{Atmosphère réelle et atmosphère standard}

L'atmosphère standard n'est qu'un modèle.

Les conditions réelles peuvent s'en écarter significativement.

Les calculateurs balistiques modernes permettent généralement de renseigner :

\begin{itemize}
\item la température ;
\item la pression ;
\item l'humidité ;
\item l'altitude.
\end{itemize}

Ils utilisent alors ces données pour corriger les calculs.

\section{Application aux essais}

Lors d'une campagne d'essais, les conditions atmosphériques doivent être
enregistrées avec soin.

Les paramètres les plus fréquemment relevés sont :

\begin{itemize}
\item la température ;
\item la pression atmosphérique ;
\item l'humidité ;
\item la vitesse du vent ;
\item la direction du vent.
\end{itemize}

Ces informations permettent :

\begin{itemize}
\item de reproduire les essais ;
\item d'expliquer certaines dispersions ;
\item de comparer les résultats à des simulations.
\end{itemize}

\begin{simulation}

La plupart des moteurs balistiques utilisent par défaut une atmosphère
standard.

Les modèles les plus avancés permettent de modifier les paramètres
atmosphériques afin de reproduire plus fidèlement les conditions réelles.

\end{simulation}

\section{À retenir}

\begin{aretenir}

\begin{itemize}
\item L'atmosphère standard est un modèle de référence.
\item Le modèle ISA est largement utilisé en aéronautique et en balistique.
\item La pression et la densité diminuent généralement avec l'altitude.
\item Les conditions atmosphériques influencent directement les trajectoires.
\item Les essais balistiques doivent toujours être associés à des mesures météorologiques.
\end{itemize}

\end{aretenir}

